\documentclass[oneside]{scrreprt}
\usepackage{arbeitsblatt}
\usepackage[bottom=2cm]{geometry}
\title{Wichtige Befehle}

\ihead{Gitkurs}
\ohead{}

\begin{document}
\section*{Unter der Haube}

\verb+git ls-files+ \\
Listet alle Dateien auf, die im Stage liegen. \\
Das schließt auch Dateien ein, die bereits committet sind. \\ 
Liefert Infos der Art\\
\verb+    100644 78981922613b2afb6025042ff6bd878ac1994e85 0	a.txt+
~\\ 
 
 
\verb+git cat-file -t <hash>+\\
Typ des Objekts wird geliefert (tree, blob, commit)
~\\

\verb+git cat-file -p <hash>+\\
Liefert den Inhalt des Objekts. Bei \emph{blob} ist es der Inhalt, bei \emph{tree} sind die Dateien im Tree.
Bei einem Commit erhält man die Message, Parent, ... und auch den Hash des trees! 

 
\verb+git rm --cached <datei>+\\
Löscht die Datei aus dem Stage. Sinnvoll, wenn die Datei in
zukünftigen Commits nicht mehr auftauchen soll. Sie muss dafür 
in die \verb+.gitignore+ aufgenommen werden.
~\\

\section*{Reparatur}
\subsection*{Reset}
\verb+git reset+\\
HEAD Entfernt eine Datei aus dem Stage, als wäre kein \verb+.add+ erfolgt. Ist quasi „einen Schritt zurück“.
Änderungen nach dem \verb+.add+, die an anderen Dateien erfolgt sind, bleiben erhalten!
~\\

\verb+git reset --hard <hash>+\\
Alles zurück auf den Stand nach dem entsprechenden Commit.
~\\

\verb+git reset --soft <hash>+\\
Der HEAD wird auf <hash> gesetzt, die Dateien in workdir und stage bleiben erhalten.
~\\

\verb+git reset --mixed <hash>+\\
Ist Standard, auch ohne mixed\\
Der HEAD wird auf <hash> gesetzt, Stage wird verworfen, workdir bleibt erhalten.
~\\



\subsection*{Restore}

\verb+git restore <datei>+\\
Holt die Version aus dem Stage und überschreibt die 
Version im Workindir.
~\\

\verb+git restore <datei> --staged+\\
Wenn im Workingdir und im Stage verschiedene Versionen
existieren, wird die aus dem Stage einfach gelöscht.\\

Existiert aktuell keine Version mehr im Workingdir, 
so wird die Datei im Stage nicht gelöscht, sondern nur
ins Workingdir verschoben.
~\\


\verb+git restore --source=<HASH> <datei>+\\
Der Befehl holt eine Datei aus einem bestimmten Commit.
Früher war das\\
 \texttt{git checkout <HASH> -\,- <datei>}\\
 Man kann hier auch noch andere Quellen angeben!


\section*{Werkzeuge}

\subsection*{Diff}

\begin{tabular}{ll}
    \verb+git diff+ & work $\leftrightarrow$ stage\\
    \verb+git diff --staged+ & stage $\leftrightarrow$ commit\\
    \verb+git diff branch1 branch2+ & branch1 $\leftrightarrow$ branch2\\
    \verb+git diff branch1...branch2+ & branch2 $\leftrightarrow$ Merge-base von branch1 und branch2\\
\end{tabular}
~\\

\subsection*{Show}
\verb+git show <HASH>+\\
Zeigt den Inhalt des Commits

\verb+git show <HASH>:datei.txt+\\
Zeigt den Inhalt der Datei zum damaligen Zeitpunkt an.

\subsection*{Log}
\verb+git log branch_1 branch_2+\\
Alle Commits in beiden Branches

\verb+git log branch_1..branch_2+\\
Alle Commits die in branch\_2, aber nicht in branch\_1 sind.

\verb+git log branch_1...branch_2+\\
Alle Commits die in genau einem Branch sind.


\end{document}